\documentclass[10pt]{article}
\linespread{1} 
\usepackage{extsizes}
\usepackage{amsmath}
\usepackage{amssymb}
\usepackage{amsthm}
\usepackage{ mathrsfs }
\usepackage{ mathtools}
\usepackage{enumerate}
\usepackage{bbm}
\usepackage{lipsum}
\usepackage{fancyhdr}
\usepackage{tikz-cd} 
\usepackage{adjustbox} 
\usepackage{ stmaryrd }
\numberwithin{equation}{section}
\theoremstyle{definition}
\usetikzlibrary{arrows}
\newcommand{\midarrow}{\tikz \draw[-triangle 90] (0,0) -- +(.1,0);}
\tikzset{commutative diagrams/.cd,
mysymbol/.style={start anchor=center,end anchor=center,draw=none}
}
\newcommand\comsymb[2][\circlearrowleft]{%
  \arrow[mysymbol]{#2}[description]{#1}}

\newtheorem{theorem}{Theorem} [section] 
\newtheorem{proposition}{Proposition}[section] 
\newtheorem{definition}{Definition}[section] 
\newtheorem{lemma}{Lemma}[section] 
\newtheorem{notation}{Notation}[section] 
\newtheorem{remark}{Remark}[section] 
\newtheorem{corollary}{Corollary} [section] 
\newtheorem{terminology}{Terminology}[section] 
\newtheorem{fact}{Fact}[section] 
\newtheorem{example}{Example}[section] 
\newtheorem{claim}{Claim}
\newenvironment{claimproof}[1]{\par\noindent\underline{Proof:}\space#1}{\hfill $\blacksquare$}

\DeclareMathOperator{\tmod}{mod}
\DeclareMathOperator{\triv}{triv}
\DeclareMathOperator{\ind}{ind}
\DeclareMathOperator{\straw}{s.t.}
\DeclareMathOperator{\ev}{ev}
\DeclareMathOperator{\pr}{pr}
\DeclareMathOperator{\Aut}{Aut}
\DeclareMathOperator{\Map}{Map}
\DeclareMathOperator{\Stab}{Stab}
\DeclareMathOperator{\Ind}{Ind}
\DeclareMathOperator{\GL}{GL}
\DeclareMathOperator{\SL}{SL}
\DeclareMathOperator{\SO}{SO}
\DeclareMathOperator{\Sp}{Sp}
\DeclareMathOperator{\esssup}{esssup}
\DeclareMathOperator{\abs}{abs}
\DeclareMathOperator{\rel}{rel}
\DeclareMathOperator{\diam}{diam}
\DeclareMathOperator{\rank}{rank}
\DeclareMathOperator{\Hom}{Hom}
\DeclareMathOperator{\Homk}{Hom_{k-alg}}
\DeclareMathOperator{\Ker}{Ker}
\DeclareMathOperator{\Image}{Im}
\DeclareMathOperator{\Dom}{Dom}
\DeclareMathOperator{\grad}{grad}
\DeclareMathOperator{\divergence}{div}
\DeclareMathOperator{\rk}{rk}
\DeclareMathOperator{\Span}{Span}
\DeclareMathOperator{\mSpec}{mSpec}
\DeclareMathOperator{\MaxSpec}{MaxSpec}
\DeclareMathOperator{\interior}{int}
\DeclareMathOperator{\supp}{supp}
\DeclareMathOperator{\id}{id}
\DeclareMathOperator{\sgn}{sgn}
\DeclareMathOperator{\Mat}{Mat}
\DeclareMathOperator{\PD}{PD}
\DeclareMathOperator{\ext}{ext}
\DeclareMathOperator{\vol}{vol}
\DeclareMathOperator{\inte}{int}
\DeclareMathOperator{\diverg}{div}
\DeclareMathOperator{\curl}{curl}
\DeclareMathOperator{\image}{im}
\DeclareMathOperator{\dR}{dR}
\DeclareMathOperator{\Ob}{Ob}
\DeclareMathOperator{\Mnfd}{Mnfd}
\DeclareMathOperator{\OpEmb}{OpEmb}
\DeclareMathOperator{\Spec}{Spec}
\DeclareMathOperator{\Spa}{Spa}
\DeclareMathOperator{\Lim}{Lim}
\DeclareMathOperator{\Alg}{\mathcal{A}}
\DeclareMathOperator{\glim}{glim}
\DeclareMathOperator{\con}{con}
\DeclareMathOperator{\Spv}{Spv}
\DeclareMathOperator{\clet}{c}
\DeclareMathOperator{\tat}{at}
\DeclareMathOperator{\Sheaf}{Sheaf}
\DeclareMathOperator{\Cont}{Cont}
%\newcommand*{\name}[\num_arguments][default values]{{\color{#1}\Large #2}}
\newcommand{\defeq}{\vcentcolon=}
\newcommand{\Ouv}{\mathscr{O}}
\newcommand{\norm}[1]{\Vert #1 \Vert}
\newcommand{\opNorm}[2]{\norm{#1}_{#2\to#2}}
\newcommand{\normL}[3]{\norm{#1}_{L^{#2}(#3)}}
\newcommand{\N}[0]{\mathbb{N}}
\newcommand{\m}[0]{\mathfrak{m}}
\newcommand{\R}[0]{\mathbb{R}}
\newcommand{\Z}[0]{\mathbb{Z}}
\newcommand{\C}[0]{\mathbb{C}}
\newcommand{\Q}[0]{\mathbb{Q}}
\newcommand{\bF}[0]{\mathbb{F}}
\newcommand{\U}[0]{\mathcal{U}}
\newcommand{\B}[0]{\mathcal{B}}
\newcommand{\G}[0]{\mathbb{G}}
\newcommand{\A}[0]{\mathbb{A}}
\newcommand{\T}[0]{\mathbb{T}}
\newcommand{\K}[0]{\mathcal{K}}
\newcommand{\Para}[0]{\mathbb{P}}
\newcommand{\M}[0]{\mathcal{M}}
\newcommand{\maniN}[0]{\mathcal{N}}
\newcommand{\cinf}[0]{\mathcal{C}^\infty}
\newcommand{\torus}[0]{\mathbb{T}}
\newcommand{\cR}[0]{\mathcal{R}}
\newcommand{\sep}[0]{\:|\:}
\newcommand{\pd}[2]{{\frac {\partial #1} {\partial #2}}}
\newcommand{\compinc}[0]{\subset \subset}
\newcommand{\sH}[0]{\mathscr{H}}
\newcommand{\Nfil}[0]{\mathcal{N}}
\newcommand{\g}[0]{\mathfrak{g}}
\newcommand{\p}[0]{\mathfrak{p}}
\newcommand{\QC}[0]{\mathfrak{Q}_C}
\newcommand{\ab}[1]{\vert #1 \vert}
\newcommand{\st}[0]{\:\straw\:}
\newcommand{\cGamma}[0]{c\Gamma}
\newcommand{\sC}[0]{\mathscr{C}}
\newcommand{\Grp}[0]{\mathrm{Grp}}
\newcommand{\op}[0]{\mathrm{op}}
\newcommand{\Func}[0]{\mathrm{Func}}
\newcommand{\sU}[0]{\mathscr{U}}
\newcommand{\cV}[0]{\mathcal{V}}
\newcommand{\sA}[0]{\mathscr{A}}
\newcommand{\sB}[0]{\mathscr{B}}
\newcommand{\sX}[0]{\mathscr{X}}
\newcommand{\sS}[0]{\mathscr{S}}
\newcommand{\sT}[0]{\mathscr{T}}
\newcommand{\cB}[0]{\mathcal{B}}
\newcommand{\cA}[0]{\mathcal{A}}
\newcommand{\cC}[0]{\mathcal{C}}
\newcommand{\cO}[0]{\mathcal{O}}
\newcommand{\cT}[0]{\mathcal{T}}
\newcommand{\cF}[0]{\mathcal{F}}
\newcommand{\cG}[0]{\mathcal{G}}
\newcommand{\nat}[0]{\mathrm{nat}}
\newcommand{\tr}[0]{\mathrm{tr}}
\newcommand{\tand}[0]{\text{ and }}
\newcommand{\tor}[0]{\text{ or }}
\newcommand{\dcirc}[0]{{\circ\circ}}
\newcommand{\otherwise}[0]{\text{ otherwise }}
\newcommand{\Rat}[0]{\mathfrak{Rat}}
\newcommand{\polring}[1]{[X_1,\cdots,X_{#1}]}
\newcommand{\compring}[1]{\langle X_1,\cdots,X_{#1}\rangle}
\newcommand{\fib}[1]{%
  \mathbin{\mathop{\times}\limits_{#1}}%
}
\newcommand{\tens}[1]{%
  \mathbin{\mathop{\otimes}\displaylimits_{#1}}%
}
\newcommand{\comppol}[1]{\langle X_1,\cdots, X_{#1}\rangle}

\title{Compact Lie Groups}
\author{Based on Lectures given by\\\\Prof. Dr. Tasho Kaletha}
\date{WSSe 24/25, University of Bonn}


\begin{document}
\maketitle

\section{Basic Lie Theory}

\subsection{Lie Groups}

\begin{definition}
    \label{definition_1_1}
    A Lie group is a pair $(G,m)$ such that $G$ is a smooth manifold and $m:G\times G\to G$ is a smooth map giving $G$ a group structure.
\end{definition}
\begin{remark}
    We can derive the smoothness of the inversion $G\ni g\stackrel{i}{\mapsto} g^{-1}\in G$ from Definition \ref{definition_1_1}.
\end{remark}

\begin{definition}
    Let $f:M\to N$ be a smooth map of smooth manifolds. $p\in M$ is said to be regular if $T_pf$ is surjective. A point $q\in N$ is a regular value if $f^{-1}(q)$ is non-empty set consisting of regular points.
\end{definition}

\begin{definition}
    A subset $N\subseteq M$ of a smooth manifold is called an embedded submanifold if for any point $p\in N$, there is a chart $(U,\varphi)$ of $M$ such that $\varphi(U\cap N)=\varphi(U)\cap (\R^k\times\{0\})$ for some $k$.
\end{definition}

\begin{theorem}
    If $q\in N$ is a regular value of $f:M\to N$ then $f^{-1}(q)$ is a closed embedded submanifold of $M$.
\end{theorem}

\begin{example}
    $(S^1,+)$ is a compact Lie group. Indeed consider 
    \begin{equation*}
    \R^2\ni(x,y)\stackrel{f}{\mapsto}x^2+y^2\in\R.
    \end{equation*}
    Then for any $(x,y)\in f^{-1}(1)$, $T_{(x,y)}f = (2x,2y)$. Hence $1$ is a regular value of $f$ and in particular, $f^{-1}(1) = S^1$ is a closed embedded submanifold.
\end{example}

\begin{lemma}
    A map $\det:\R^{n\times n}\to \R$ is differentiable at any $A\in\GL_n(\R)$. Furthermore, the derivative is given explicitly as follows.
    \begin{equation*}
        d[\det](A)(H) = \det(A)\tr(A^{-1}H).
    \end{equation*}
\end{lemma}

\begin{proof}$\:$
    \par\textbf{Base Case : $A=I$}. 
    \begin{align*}
        d[\det](I)(H) &= \det(I+H)-\det(I)\mod{o(H)}.
    \end{align*}
    Write $H = (h_1|\cdots|h_n)$ where $h_1,\cdots,h_n$ are column vectors. Then 
    \begin{align*}
        \det(I+H) &= \det(I)+\sum_{i=1}^n\det(e_1|\cdots|h_i|\cdots|e_n)\mod{o(H)}\\
        & = \det(I) +\tr(H).
    \end{align*}
    Hence \textbf{Base Case} is shown.
    \par\textbf{General Case}. Observe 
    \begin{equation}
        \det(A+H)-\det(A) = \det(A)(\det(I+A^{-1}H)-\det(I)).
    \end{equation}
\end{proof}

\begin{corollary}
    $1$ is a regular value of $\det$.
\end{corollary}
\begin{proof}
    Follows from the non-degeneracy of $\tr$ (i.e., $\forall Y\in\Mat_{n\times n},\tr(XY) = 0\Rightarrow X=0$).
\end{proof}

\begin{corollary}
    $\SL_n(\R)$ is a closed embedded submanifold of $\GL_n(\R)$.
\end{corollary}

\begin{proposition}
    For a smooth map $f:M\to N$ and $q\in N$ a regular value, we have 
    \begin{equation*}
        T_pf^{-1}(q) = d_pf^{-1}(0)\subseteq T_pM.
    \end{equation*}
\end{proposition}

\begin{definition}
    \begin{equation*}
        \mathfrak{sl}_n \defeq \tr^{-1}(0).
    \end{equation*}
\end{definition}

\begin{example}
    Let $\langle\cdot,\cdot\rangle:V\times V\to \R$ be a non-degenerate bilinear form. The orthogonal group with respect to $\langle\cdot,\cdot\rangle$ is  
    \begin{equation*}
        O(\langle\cdot,\cdot\rangle) \defeq \{A\in\GL_n(\R)\sep \forall x,y, \langle Ax,Ay\rangle = \langle x,y\rangle\}.
    \end{equation*}
\end{example}

\begin{notation}
Due to Sylvester's theorem, there is an isomorphism between $\langle\cdot,\cdot\rangle$ and the quadratic form given by the matrix
\begin{equation*}
    \mathrm{diag}(\underbrace{1,\cdots,1}_{p},\underbrace{-1,\cdots,-1}_{q}).
\end{equation*}
With this identification, we may write $O(\langle\cdot,\cdot\rangle) = O(p,q)$.
\end{notation}
\begin{remark}
    $O(p,q)$ is compact if and only if $p=0$ or $q=0$.
\end{remark}

\begin{proposition}
    For the standard inner product $\langle\cdot,\cdot\rangle$ in $\R^n$, the map
    \begin{equation*}
        \Mat_{n\times n}(\R)\ni A\mapsto A^*\langle\cdot|\cdot\rangle \defeq \langle A\cdot|A\cdot\rangle\in S^2(\R^n)^*
    \end{equation*}
    is differentiable everywhere and submersive at any $A\in O(p,q)$.
\end{proposition}
\begin{proof}
    \begin{align*}
        (A+H)^*\langle x|y\rangle-A^*\langle x|y\rangle & \equiv \langle Hx|Ay\rangle+\langle Ax|Hy\rangle +\langle Hx|Hy\rangle\mod{o(H)}\\
        &\equiv \langle Hx|Ay\rangle+\langle Ax|Hy\rangle\mod{o(H)}.
    \end{align*}
    \par For the surjectivity, we first prove for $A=I$. Since any such symmetric bilinear form is given by $B\in\Mat_{n\times n}(\R)$ where $B$ is a symmetric matrix and $\langle Bx|y\rangle$. Using the symmetry of $B$, we see that by setting $H={\frac B 2}$, the surjectivity is shown.
    \par For the general case, set $H' = AH$, then $(H') = H^*A^* = HA^{-1}$ hence the base case can be applied.
\end{proof}



\begin{thebibliography}{9}
\bibitem{Borel} Linear algebraic groups.
\bibitem{Milne} Linear algebraic groups.
\end{thebibliography}
\end{document}